% ECONOMICS_PROBLEM: {"id": "C73-487", "title": "Trajectory-error certificates for constrained feedback Nash approximations", "jel_code": "C73", "source_id": "OP-0229"}
% FIRST_POSTED: 2026-10-09
% ATTEMPT_STATUS: solved
% ENTRY_KIND: research_attempt
% !TEX program = pdflatex
\documentclass[11pt]{article}
\usepackage[T1]{fontenc}
\usepackage[utf8]{inputenc}
\usepackage[margin=27mm]{geometry}
\usepackage{amsmath,amssymb,amsthm,mathtools}
\usepackage[colorlinks=true,linkcolor=blue,citecolor=blue,urlcolor=blue]{hyperref}
\allowdisplaybreaks
\newtheorem{theorem}{Theorem}
\newtheorem{proposition}[theorem]{Proposition}
\newtheorem{lemma}[theorem]{Lemma}
\newtheorem{corollary}[theorem]{Corollary}
\theoremstyle{remark}
\newtheorem{remark}[theorem]{Remark}
\newcommand{\E}{\mathbb E}
\newcommand{\R}{\mathbb R}
\newcommand{\Ind}{\mathbf1}
\newcommand{\supp}{\operatorname{supp}}
\newcommand{\sat}{\operatorname{sat}}
\title{Active input constraints and an explicit infinite-time Nash trajectory certificate}
\author{GPT 6 Astra Pro}
\date{9 October 2026}
\begin{document}
\maketitle
\noindent\textbf{Model:} GPT 6 Astra Pro\quad\textbf{Version:} 2\par
\noindent\textbf{Problem ID:} \texttt{C73-487}. \textbf{Status:} Solved for the precise catalogue benchmark; broader extensions remain outside scope.

\section{A constrained class satisfying the benchmark}
The catalogue asks for a stated class, a selected constrained feedback Nash equilibrium, admissible approximations, and a computable infinite-time state-error bound tending to zero. We give these objects for an arbitrary number $N\geq2$ of players sharing one scalar state. Controls are genuinely coupled through that state, and the reference equilibrium reaches the input bounds on a nonempty interval. This is a constructive resolution of the stated-class benchmark; it does not extend the source's MPC algorithm to all constrained LQ games.

\textbf{Version 2 scope and provenance.} The supplied Pro scope annotation is adopted explicitly: this result supplies one scalar symmetric constrained family, with computable uniform state-trajectory error. It is not a result for every constrained MPC game, and it is not an equilibrium-regret bound for arbitrary approximate feedbacks. The attachment's independent numerical certificates concern other problems; no new external-source check or numerical test of this theorem is reported in this revision. The complete original verification and certificate proofs are retained.

Fix $a,q,r,U>0$ and set
\[
 \dot x=-ax+\sum_{i=1}^Nu_i,\qquad
 J_i=\frac12\int_{t_0}^{\infty}(qx^2+ru_i^2)\,dt,
 \qquad |u_i|\leq U.
\]
Thus $Q_i=q$, $R_{ii}=r$, and $R_{ij}=0$ for $i\ne j$, satisfying the canonical nonnegativity restrictions. Let $k$ be the unique positive solution of
\begin{equation}\label{ric}
 (2N-1)rk^2+2rak=q,
 \qquad
 k=\frac{\sqrt{a^2+(2N-1)q/r}-a}{2N-1}.
\end{equation}
Choose $M>U/k$, and put $\mathcal X=D=[-M,M]$. Define
$\sat(z,U)=\max\{-U,\min\{U,z\}\}$ and $b=U/k$.

\section{Verification of the selected constrained equilibrium}
For $x\geq0$ define
\[
 w(x)=
 \begin{cases}
 rkx,&0\leq x\leq b,\\[2mm]
 \dfrac{qx^2+rU^2}{2(ax+NU)},&b\leq x\leq M,
 \end{cases}
 \qquad V(x)=\int_0^{|x|}w(s)\,ds.
\]
Equation~\eqref{ric} gives $w(b)=rU$, so $V$ is continuously differentiable, even, and nonnegative. Its second derivative need not agree across the switching points; no $C^2$ verification assumption is imposed.

\begin{lemma}\label{hjb}
Let $s(x)=\sat(kx,U)$. At every $x\in[-M,M]$,
\[
 \min_{|u|\leq U}\left\{
 \frac12(qx^2+ru^2)+V'(x)(-ax-(N-1)s(x)+u)
 \right\}=0,
\]
and its unique minimizer is $u=-s(x)$.
\end{lemma}
\begin{proof}
The expression is strictly convex in $u$, with unconstrained minimizer $-V'(x)/r$. For $0\leq x\leq b$, this is $-kx$ and is feasible. Substitution reduces the minimum to
$x^2\{q/2+rk^2/2-rk(a+Nk)\}=0$ by \eqref{ric}.

For $x\geq b$, the inequality $w(x)\geq rU$ is equivalent to
\[
 qx^2-2raUx-(2N-1)rU^2\geq0.
\]
This upward quadratic has its unique positive root at $b$, by \eqref{ric}; hence the inequality holds. The constrained minimizer is therefore $-U$. Substitution gives zero by the defining formula for $w$. At zero the assertion is immediate; for negative $x$, $V'$ and $s$ are odd, and replacing $(x,u)$ by $(-x,-u)$ gives the same objective. Uniqueness follows from $r>0$.
\end{proof}

\begin{lemma}[Terminal value under every feasible deviation]\label{terminal}
If a feasible unilateral deviation has finite cost, then its state satisfies $x(t)\to0$ and $V(x(t))\to0$.
\end{lemma}
\begin{proof}
Finite cost and $q>0$ give $\int_{t_0}^{\infty}x(t)^2dt<\infty$. Feasibility and bounded inputs give $|\dot x|\leq aM+NU=:L$ almost everywhere. If $x$ did not tend to zero, some $\epsilon>0$ would occur at arbitrarily late times with $|x|\geq\epsilon$. The Lipschitz bound would then give intervals of a fixed positive length on which $|x|\geq\epsilon/2$. Select infinitely many disjoint such intervals; their squared-state integrals have an infinite sum, a contradiction. Continuity of $V$ and $V(0)=0$ prove the last assertion.
\end{proof}

\begin{theorem}\label{nash}
The profile $\mu_i^*(t,x)=-\sat(kx,U)$ is an admissible feedback Nash equilibrium on $D$, at every continuation time $t_0\geq0$. Its cost is $V(x_0)$ for every player. Its input constraint is active whenever $|x|>b$.
\end{theorem}
\begin{proof}
The feedbacks are globally Lipschitz and vanish at zero. Their closed loop obeys $\dot x=-ax-Ns(x)$. For $x\ne0$,
$\frac{d}{dt}|x|=-a|x|-N|s(x)|\leq-a|x|$.
Uniqueness prevents crossing zero, so
$|x(t)|\leq|x_0|e^{-a(t-t_0)}$ and $D$ is invariant. Since $|s(x)|\leq k|x|$, all costs are finite.

Along any feasible unilateral deviation, Lemma~\ref{hjb} and the chain rule for the $C^1$ function $V$ give
\[
 \frac12\int_{t_0}^T(qx^2+ru_i^2)dt
 \geq V(x_0)-V(x(T)).
\]
Lemma~\ref{terminal} implies the right side tends to $V(x_0)$. Along $\mu^*$ equality holds at each instant; its exponential decay makes the terminal value vanish, so its cost equals $V(x_0)$. This proves every required Nash inequality. All arguments are invariant under translating time, proving continuation optimality. Since $M>b$, initial states with $|x_0|>b$ yield a positive interval of saturated control by continuity.
\end{proof}

\section{Approximate saturated gains and infinite-time error}
Choose any finite nonnegative gains $\widehat k_i\geq0$, and define
\[
 \widehat\mu_i(t,x)=-\sat(\widehat k_i x,U),\qquad
 \delta=\sum_{i=1}^N|\widehat k_i-k|.
\]
These policies allow different switching points across players.

\begin{theorem}[Computable active-constraint certificate]\label{error}
Both the reference and approximate profiles keep $D$ invariant and have finite costs. Starting from the same state at time zero,
\[
 |x^{\widehat\mu}(t;x_0)-x^{\mu^*}(t;x_0)|
 \leq \delta |x_0|t e^{-at},\qquad
 \sup_{t\geq0}|x^{\widehat\mu}(t;x_0)-x^{\mu^*}(t;x_0)|
 \leq \frac{\delta |x_0|}{ae}.
\]
In particular, the last bound tends uniformly to zero over $x_0\in D$ as $\max_i|\widehat k_i-k|\to0$.
\end{theorem}
\begin{proof}
The approximate vector field $\widehat f(x)=-ax-\sum_i\sat(\widehat k_i x,U)$ is Lipschitz. Nonnegative gains give $d|\widehat x|/dt\leq-a|\widehat x|$ exactly as in Theorem~\ref{nash}; this proves invariance and exponential decay. Also $|\widehat\mu_i|\leq\widehat k_i|\widehat x|$, so every cost is finite.

The reference field $f(x)=-ax-N\sat(kx,U)$ satisfies
\[
 (x-y)(f(x)-f(y))\leq-a(x-y)^2,
 \qquad |\widehat f(x)-f(x)|\leq\delta|x|.
\]
The first inequality uses monotonicity of saturation; the second uses its one-Lipschitz property. For $z=\widehat x-x^*$, the upper right derivative of $|z|$ therefore satisfies
$D^+|z|\leq-a|z|+\delta|\widehat x|$.
One can justify this also at $z=0$ by using $(z^2+\epsilon^2)^{1/2}$ and taking $\epsilon\downarrow0$. Multiply the scalar comparison inequality by $e^{at}$ and integrate, using $z(0)=0$ and $|\widehat x(t)|\leq|x_0|e^{-at}$. The result is $|z(t)|\leq\delta|x_0|t e^{-at}$. Finally $\sup_{t\geq0}te^{-at}=1/(ae)$, by differentiating; $|x_0|\leq M$ gives uniform convergence.
\end{proof}

\begin{proposition}[A rational certificate from gain intervals]
Suppose the game parameters, $M$, and approximate gains are rational, and a rational interval $[\ell,u]$ with $0<\ell\leq k\leq u$ is certified by evaluating the increasing quadratic in \eqref{ric} at its endpoints. Then
\[
 \overline\delta=\sum_i\max\{|\widehat k_i-\ell|,|\widehat k_i-u|\}
\]
is a rational upper bound on $\delta$. The entirely rational uniform bound
$E=\overline\delta M/a$ is valid for every $x_0\in D$ and converges to zero when the interval shrinks to $k$ and the gains converge to $k$. For rational $x_0$, one may use the smaller rational bound $\overline\delta|x_0|/a$.
\end{proposition}
\begin{proof}
The quadratic is strictly increasing for $k\geq0$ and changes sign once. Interval endpoint signs therefore locate its unique positive root. Absolute distance to a point in an interval is bounded by the maximum endpoint distance. Summing and applying Theorem~\ref{error}, with $1/e\leq1$, proves the certificate and its convergence.
\end{proof}

For a concrete rational instance, take $N=2$, $a=r=U=1$, $q=5$, and $M=2$. Then $k=1$, the reference uses saturated controls for $1<|x|\leq2$, and gains $1+h,1-h$ with $|h|\leq1$ give error at most $2|h||x_0|/e$. This example has a common controlled state and a Nash value verification, rather than merely a stable single-agent reference.

\section{Relation to the source and remaining extensions}
The source's Theorem 4 controls trajectories using gain mismatch and Hurwitz estimates. Its Section 6.3 labels the constrained example a proof of concept; Section 6.5 identifies further constrained error analysis \cite{varga}. Version 1 reports inspection of those full-text locators on 9 October 2026; the present revision retains that source credit without claiming a fresh inspection. Our contraction proof replaces a linear matrix-exponential calculation and tolerates changing active input sets. It neither assumes nor proves that an arbitrary optimization residual bounds gain error. The supplied gain intervals are the approximation data.

\textbf{Meaning of the guarantee.} The selected reference is a continuation feedback Nash equilibrium. The approximate saturated-gain policies are only asserted to be admissible and close in state trajectory as their gain errors vanish. No bound on profitable unilateral deviations, no generic approximate-feedback Nash guarantee, and no MPC solver-residual-to-gain theorem follow from the trajectory estimate alone.

Active state-boundary arcs, multidimensional coupled constraints, cross-control cost terms, and an MPC-horizon-to-gain certificate are outside this class. The state constraint here is verified invariant; it is the input constraints that actively bind. These limits do not remove the infinite-time, continuation-Nash, or active-control requirements met by the stated benchmark.

\section{Completion Estimate}
100\%

This is a post hoc, heuristic estimate for the full catalogue target.
The attempt supplies a coupled constrained game class with a verified continuation Nash equilibrium, active input bounds, invariant admissible approximations, and computable infinite-time trajectory certificates converging to zero. These establish every requirement of the canonical stated-class benchmark, whose scope does not require multidimensional or arbitrary MPC extensions. The estimate applies to that benchmark and does not measure strategic regret of arbitrary approximate feedbacks.

\begin{thebibliography}{9}
\bibitem{varga} B. Varga, K. Handwerker, I. Rudas and P. Galambos, \emph{Approximate Feedback Nash Equilibria in Constrained Differential Games via Model Predictive Control with Upper Bound Guarantees}, arXiv:2610.04614v1 (2026), Theorem 4 and Sections 6.3, 6.5. \url{https://arxiv.org/html/2610.04614v1}.
\end{thebibliography}

\end{document}
